\begin{remark*}[Exposition]
Limit points and isolated points separate two ways a subset can contain or
approach a point.  A limit point is repeatedly approached by other points of
the set at every scale.  An isolated point belongs to the set but has a small
ball in which it is the only point of the set.
\end{remark*}

\begin{definitionbox}{Definition --- Metric Limit Point and Isolated Point}
\begin{definition}[Metric Limit Point and Isolated Point]
\label{def:metric-limit-point-isolated-point}
Let \(E\) be a subset of a metric space \((X,d)\), and let \(x\in X\).
\begin{enumerate}
  \item The point \(x\) is a \emph{limit point of \(E\)} if, for every
        \(\varepsilon>0\), the ball \(B_d(x;\varepsilon)\) contains a point of
        \(E\setminus\{x\}\).
  \item The point \(x\) is an \emph{isolated point of \(E\)} if \(x\in E\) and
        \(x\) is not a limit point of \(E\).
\end{enumerate}
\end{definition}
\end{definitionbox}

\begin{remark*}[Standard quantified statement]
\[
x\text{ is a limit point of }E
\quad\Longleftrightarrow\quad
\forall\varepsilon>0{:}\;
B_d(x;\varepsilon)\cap(E\setminus\{x\})\neq\varnothing.
\]
\[
x\text{ is an isolated point of }E
\quad\Longleftrightarrow\quad
x\in E
\;\land\;
\exists\varepsilon>0{:}\;
B_d(x;\varepsilon)\cap E=\{x\}.
\]
\end{remark*}

\begin{remark*}[Interpretation]
A limit point of \(E\) may lie in \(E\) or outside \(E\); what matters is that
points of \(E\) distinct from \(x\) occur arbitrarily close to \(x\).  An
isolated point must belong to \(E\), and some positive-radius ball around it
contains no other point of \(E\).
\end{remark*}

\begin{remark*}[Notation]
The set of limit points of \(E\) is denoted by \(E'\), read as
\emph{\(E\)-prime}.
\end{remark*}

\begin{remark*}[Exposition]
Limit points are also called \emph{cluster points} or \emph{accumulation
points}.  In later work with functions on subsets of metric spaces, limits of
the form \(\lim_{y\to x}f(y)\) are only meaningful when \(x\) is a limit point
of the domain being approached.
\end{remark*}

\begin{remark*}[Examples]
In the usual metric on \(\mathbb{R}\), if
\[
  E=(0,1)\cup\{2\},
\]
then \(2\) is an isolated point of \(E\), while \([0,1]\) is the set of limit
points of \(E\).

There are also metric spaces with infinite subsets that have no limit points.
For example, if \(X\) is an infinite discrete metric space and \(E=X\), then
\(E'=\varnothing\).  Also, in \(X=\mathbb{R}\setminus\mathbb{Q}\), the set
\[
  E=\{\sqrt{2}/n:n\in\mathbb{N}\}
\]
has no limit point in \(X\).
\end{remark*}

\begin{dependencies}
\begin{itemize}
  \item \hyperref[def:metric-space]{Definition: Metric Space}
  \item \hyperref[def:metric-open-ball]{Definition: Open Ball}
\end{itemize}
\end{dependencies}

\begin{theorembox}{Theorem}
\begin{theorem}[Closed Sets Contain Their Limit Points]
\label{thm:metric-closed-sets-contain-limit-points}
Let \(X\) be a metric space and let \(E\subseteq X\).  Then \(E\) is closed if
and only if \(E'\subseteq E\).

\smallskip
\noindent
\hyperref[prf:metric-closed-sets-contain-limit-points]{\textit{Go to proof.}}
\end{theorem}
\end{theorembox}

\begin{remark*}[Standard quantified statement]
\[
E\text{ is closed in }X
\quad\Longleftrightarrow\quad
E'\subseteq E.
\]
\end{remark*}

\begin{remark*}[Interpretation]
Closed sets contain all points that can be approached by other points of the
set.  Equivalently, no limit point of a closed set is missing from the set.
\end{remark*}

\begin{remark*}[Historical note]
This result corresponds to Theorem~3.10 in Kainth's
\emph{A Comprehensive Textbook on Metric Spaces}
\cite{KainthComprehensiveMetricSpaces2023}.
\end{remark*}

\begin{dependencies}
\begin{itemize}
  \item \hyperref[def:metric-closed-set]{Definition: Metric Closed Set}
  \item \hyperref[def:metric-limit-point-isolated-point]{Definition: Metric Limit Point and Isolated Point}
\end{itemize}
\end{dependencies}

\begin{theorembox}{Theorem}
\begin{theorem}[Sequential Characterization of Limit Points]
\label{thm:metric-limit-point-sequential-characterization}
Let \((X,d)\) be a metric space, let \(E\subseteq X\), and let \(x\in X\).
Then \(x\in E'\) if and only if there exists a sequence of distinct terms in
\(E\) that converges to \(x\).

\smallskip
\noindent
\hyperref[prf:metric-limit-point-sequential-characterization]{\textit{Go to proof.}}
\end{theorem}
\end{theorembox}

\begin{remark*}[Standard quantified statement]
\[
x\in E'
\quad\Longleftrightarrow\quad
\exists \{x_n\}\subseteq E{:}\;
\bigl(\forall n\neq m,\;x_n\neq x_m\bigr)
\;\land\;
x_n\to x.
\]
\end{remark*}

\begin{remark*}[Predicate reading]
\[
\operatorname{IsCluster}(x,E,(X,d))
\Longleftrightarrow
\exists (x_n)_{n\in\mathbb{N}}{:}\;
\operatorname{ConvergesTo}(x_n,x,(X,d))
\land
\forall n\neq m,\;x_n\neq x_m
\]
means that \(x\) is a limit point of \(E\) exactly when \(E\) supplies a
sequence of distinct terms converging to \(x\).
\end{remark*}

\begin{remark*}[Interpretation]
Limit points can be detected sequentially.  A point is a limit point exactly
when the set supplies infinitely many distinct approximating terms converging
to that point.
\end{remark*}

\begin{remark*}[Historical note]
This result corresponds to Theorem~3.11 in Kainth's
\emph{A Comprehensive Textbook on Metric Spaces}
\cite{KainthComprehensiveMetricSpaces2023}.
\end{remark*}

\begin{dependencies}
\begin{itemize}
  \item \hyperref[def:metric-limit-point-isolated-point]{Definition: Metric Limit Point and Isolated Point}
  \item \hyperref[def:metric-sequential-convergence]{Definition: Metric Sequential Convergence}
\end{itemize}
\end{dependencies}

\begin{corollarybox}{Corollary}
\begin{corollary}[Neighborhood Characterization of Limit Points]
\label{cor:metric-limit-point-neighborhood-characterization}
Let \((X,d)\) be a metric space, let \(E\subseteq X\), and let \(x\in X\).
Then \(x\in E'\) if and only if every neighborhood of \(x\) contains
infinitely many points of \(E\).

\smallskip
\noindent
\hyperref[prf:metric-limit-point-neighborhood-characterization]{\textit{Go to proof.}}
\end{corollary}
\end{corollarybox}

\begin{remark*}[Standard quantified statement]
\[
x\in E'
\quad\Longleftrightarrow\quad
\forall U{:}\;
\bigl(U\text{ is a neighborhood of }x\bigr)
\Longrightarrow
U\cap E\text{ is infinite}.
\]
\end{remark*}

\begin{remark*}[Interpretation]
The neighborhood form strengthens the ball definition into a count statement:
near a limit point, every neighborhood contains not just one nearby point of
\(E\), but infinitely many.
\end{remark*}

\begin{remark*}[Historical note]
This result corresponds to Corollary~3.12 in Kainth's
\emph{A Comprehensive Textbook on Metric Spaces}
\cite{KainthComprehensiveMetricSpaces2023}.
\end{remark*}

\begin{dependencies}
\begin{itemize}
  \item \hyperref[def:metric-neighborhood]{Definition: Metric Neighborhood}
  \item \hyperref[thm:metric-limit-point-sequential-characterization]{Theorem: Sequential Characterization of Limit Points}
\end{itemize}
\end{dependencies}

\begin{theorembox}{Theorem}
\begin{theorem}[Bolzano--Weierstrass for Bounded Infinite Subsets]
\label{thm:metric-bolzano-weierstrass-bounded-infinite-subsets}
Every bounded infinite subset of \(\mathbb{R}^m\) has a limit point in
\(\mathbb{R}^m\).

\smallskip
\noindent
\hyperref[prf:metric-bolzano-weierstrass-bounded-infinite-subsets]{\textit{Go to proof.}}
\end{theorem}
\end{theorembox}

\begin{remark*}[Standard quantified statement]
\[
E\subseteq\mathbb{R}^m,\quad
E\text{ bounded},\quad
E\text{ infinite}
\quad\Longrightarrow\quad
\exists x\in\mathbb{R}^m{:}\;x\in E'.
\]
\end{remark*}

\begin{remark*}[Interpretation]
In finite-dimensional Euclidean space, boundedness prevents an infinite set
from spreading out completely.  Some point of the ambient space must be
approached by infinitely many points of the set.
\end{remark*}

\begin{remark*}[Historical note]
This result corresponds to Theorem~3.13 in Kainth's
\emph{A Comprehensive Textbook on Metric Spaces}
\cite{KainthComprehensiveMetricSpaces2023}.
\end{remark*}

\begin{dependencies}
\begin{itemize}
  \item \hyperref[thm:metric-euclidean-bolzano-weierstrass]{Theorem: Bolzano--Weierstrass in \(\mathbb{R}^m\)}
  \item \hyperref[thm:metric-limit-point-sequential-characterization]{Theorem: Sequential Characterization of Limit Points}
\end{itemize}
\end{dependencies}

\begin{propositionbox}{Proposition}
\begin{proposition}[The Derived Set is Closed]
\label{prop:metric-derived-set-closed}
The set of limit points of any set is closed.

\smallskip
\noindent
\hyperref[prf:metric-derived-set-closed]{\textit{Go to proof.}}
\end{proposition}
\end{propositionbox}

\begin{remark*}[Standard quantified statement]
\[
E\subseteq X
\quad\Longrightarrow\quad
E'\text{ is closed in }X.
\]
\end{remark*}

\begin{remark*}[Interpretation]
Taking limit points produces a closed set.  Once a point is approached by
limit points of \(E\), it is itself forced to be a limit point of \(E\).
\end{remark*}

\begin{remark*}[Historical note]
This result corresponds to Proposition~3.15 in Kainth's
\emph{A Comprehensive Textbook on Metric Spaces}
\cite{KainthComprehensiveMetricSpaces2023}.
\end{remark*}

\begin{dependencies}
\begin{itemize}
  \item \hyperref[def:metric-limit-point-isolated-point]{Definition: Metric Limit Point and Isolated Point}
  \item \hyperref[thm:metric-closed-sets-contain-limit-points]{Theorem: Closed Sets Contain Their Limit Points}
\end{itemize}
\end{dependencies}
